% Figure for Relativity §C1.4 (Model §C1.4.1, Theorem §C1.4.2): rigid rocket in units 1/g = 1, height h = 0.5 (gh = 0.5, exaggerated). Floor X = cosh(eta), top X = 1.5 cosh(eta), X = x + 1/g, t = rho sinh(eta). Pulses emitted at eta = 0 and 0.4 arrive at eta + ln 1.5. Dots every 0.2 of eta: floor proper time 0.2, top 0.3. Computed in Python (exact hyperbolic motion).
\begin{tikzpicture}[>={Stealth[length=2.0mm]}, font=\small, x=2.6cm, y=2.6cm]
  \draw[->, black!55] (-1.15,0) -- (2.45,0) node[right] {$x$};
  \draw[->, black!55] (0,-0.08) -- (0,2.3) node[above] {$t$};
  \foreach \k in {1,2} \draw[black!55] (\k,0.02) -- (\k,-0.02) node[below, font=\scriptsize] {$\k/g$};
  \foreach \k in {1,2} \draw[black!55] (0.02,\k) -- (-0.02,\k) node[left, font=\scriptsize] {$\k/g$};
  \fill[black] (-1,0) circle (1.6pt) node[below, font=\scriptsize, align=center] {pivot\\$x = -1/g$};
  \draw[black!45, dashed] (-1,0) -- (0.8919,0.7188) node[above right, font=\scriptsize, black!60, align=left] {constant $\eta$:\\simultaneous in the rocket};
  \draw[figB, very thick] plot coordinates {(0.0000,0.0000) (0.0003,0.0246) (0.0012,0.0492) (0.0027,0.0738) (0.0048,0.0985) (0.0076,0.1232) (0.0109,0.1480) (0.0148,0.1729) (0.0194,0.1979) (0.0246,0.2230) (0.0304,0.2482) (0.0368,0.2736) (0.0438,0.2992) (0.0515,0.3250) (0.0598,0.3509) (0.0687,0.3771) (0.0783,0.4034) (0.0886,0.4301) (0.0995,0.4569) (0.1110,0.4841) (0.1233,0.5116) (0.1362,0.5393) (0.1498,0.5674) (0.1641,0.5958) (0.1791,0.6246) (0.1948,0.6538) (0.2112,0.6834) (0.2284,0.7133) (0.2463,0.7437) (0.2649,0.7746) (0.2843,0.8059) (0.3045,0.8377) (0.3255,0.8701) (0.3473,0.9029) (0.3699,0.9363) (0.3933,0.9702) (0.4176,1.0048) (0.4427,1.0399) (0.4687,1.0757) (0.4956,1.1121) (0.5234,1.1492) (0.5521,1.1870) (0.5817,1.2255) (0.6123,1.2648) (0.6439,1.3048) (0.6765,1.3456) (0.7101,1.3872) (0.7447,1.4296) (0.7803,1.4729) (0.8171,1.5171) (0.8549,1.5623) (0.8939,1.6083) (0.9340,1.6554) (0.9752,1.7034) (1.0177,1.7525) (1.0614,1.8026) (1.1063,1.8538) (1.1525,1.9061) (1.2000,1.9596) (1.2488,2.0143)};
  \draw[figR, very thick] plot coordinates {(0.5000,0.0000) (0.5003,0.0287) (0.5011,0.0575) (0.5025,0.0862) (0.5044,0.1150) (0.5069,0.1439) (0.5099,0.1728) (0.5135,0.2017) (0.5176,0.2307) (0.5223,0.2598) (0.5276,0.2890) (0.5334,0.3184) (0.5398,0.3478) (0.5467,0.3773) (0.5542,0.4070) (0.5623,0.4369) (0.5710,0.4669) (0.5802,0.4971) (0.5900,0.5274) (0.6004,0.5580) (0.6114,0.5887) (0.6230,0.6197) (0.6351,0.6509) (0.6479,0.6823) (0.6613,0.7140) (0.6753,0.7460) (0.6899,0.7782) (0.7051,0.8107) (0.7209,0.8435) (0.7374,0.8766) (0.7545,0.9101) (0.7722,0.9438) (0.7906,0.9780) (0.8097,1.0124) (0.8294,1.0473) (0.8498,1.0825) (0.8709,1.1182) (0.8927,1.1542) (0.9151,1.1907) (0.9383,1.2276) (0.9621,1.2649) (0.9867,1.3027) (1.0120,1.3410) (1.0381,1.3798) (1.0649,1.4191) (1.0925,1.4589) (1.1208,1.4992) (1.1499,1.5401) (1.1798,1.5816) (1.2105,1.6236) (1.2420,1.6663) (1.2743,1.7095) (1.3075,1.7534) (1.3415,1.7979) (1.3763,1.8431) (1.4121,1.8890) (1.4487,1.9355) (1.4862,1.9828) (1.5247,2.0307) (1.5640,2.0795)};
  \fill[figB] (0.0201,0.2013) circle (1.3pt);
  \fill[figB] (0.0811,0.4108) circle (1.3pt);
  \fill[figB] (0.1855,0.6367) circle (1.3pt);
  \fill[figB] (0.3374,0.8881) circle (1.3pt);
  \fill[figB] (0.5431,1.1752) circle (1.3pt);
  \fill[figB] (0.8107,1.5095) circle (1.3pt);
  \fill[figB] (1.1509,1.9043) circle (1.3pt);
  \fill[figR] (0.5301,0.3020) circle (1.3pt);
  \fill[figR] (0.6216,0.6161) circle (1.3pt);
  \fill[figR] (0.7782,0.9550) circle (1.3pt);
  \fill[figR] (1.0062,1.3322) circle (1.3pt);
  \fill[figR] (1.3146,1.7628) circle (1.3pt);
  \draw[figO, thick, ->] (0.0000,0.0000) -- (0.6250,0.6250);
  \fill[figO] (0.0000,0.0000) circle (1.8pt); \fill[figO] (0.6250,0.6250) circle (1.8pt);
  \draw[black, thick, ->] (0.6250,0.6250) -- (0.7481,0.9450);
  \draw[figO, thick, ->] (0.0811,0.4108) -- (1.0135,1.3431);
  \fill[figO] (0.0811,0.4108) circle (1.8pt); \fill[figO] (1.0135,1.3431) circle (1.8pt);
  \draw[black, thick, ->] (1.0135,1.3431) -- (1.2269,1.6631);
  \node[figB, font=\scriptsize, anchor=east] at (0.0311,0.2054) {$\Delta\tau_E = 0.4/g$};
  \draw[figB, thin] (-0.0200,0.0000) -- (0.0511,0.4108);
  \node[figR, font=\scriptsize, anchor=west] at (1.1335,1.0341) {$\Delta\tau_R = 0.6/g = (1 + gh)\,\Delta\tau_E$};
  \node[figB, font=\scriptsize, anchor=east] at (0.62,1.72) {floor $E$: proper acceleration $g$};
  \node[figR, font=\scriptsize, anchor=south west] at (1.3640,2.0795) {top $R$: proper acceleration $g/(1+gh)$};
  \node[figO, font=\scriptsize, anchor=north west] at (0.3425,0.2925) {light};
  \node[font=\scriptsize, anchor=west, align=left] at (0.95,0.25) {$\nu_R/\nu_E = 1/(1+gh) = 2/3$\\$gh = 0.5$ (exaggerated)};
\end{tikzpicture}
